\documentclass[10pt]{article}

% graphicx package, useful for including eps, pdf, and raster graphics
\usepackage{graphicx}
\DeclareGraphicsExtensions{.png,.png,.jpg,.pdf}

% basic packages
\usepackage{fancyhdr}
\usepackage{color}
\usepackage{parskip}
\usepackage{float}
\usepackage{microtype}
\usepackage{url}
\usepackage{hyperref}
\usepackage{booktabs}
\usepackage{makecell}
\usepackage{multirow}
\usepackage{pbox}
\usepackage{enumitem}

% fonts
\usepackage{fontspec}
\setmainfont{DroidSans}[
  Extension   = .ttf,
  UprightFont = *,
  BoldFont    = *-Bold ]
\setmonofont{DroidSansMono}[
  Extension   = .ttf,
  UprightFont = * ]

\hypersetup{colorlinks,urlcolor=blue}
\urlstyle{same}

% page layout
\usepackage[top=1in, bottom=1in, left=0.75in, right=0.75in, includefoot, heightrounded]{geometry}

% Useful for seeing the layout as you tweak it.
%\usepackage{showframe}

\setlength{\headheight}{1.1in}
\pagestyle{fancy}
\fancyhf{}

\lhead{\includegraphics[height=1in,clip]{logo}}
\rhead{\pbox[b][1in][c]{\textwidth}{
    \small
    \textbf{Test Performed By}\\
    Northwest Genomics Center\\
    University of Washington School of Medicine\\
    Genome Sciences Box 355065\\
    3720 15th Ave NE\\
    Seattle, WA 98195-5065}}

\newcommand{\PageLine}{\rule{\textwidth}{0.25mm}}

\begin{document}

\begin{description}[font=\normalfont,align=left,labelwidth=12em]
\item [\textbf{\VAR{pat_name|e}}]
\item [\textbf{Date of Birth}] \VAR{birth_date|e}
\item [\textbf{Specimen}]
  %- if swab_type == "ans"
  Anterior Nasal,
  %- elif swab_type == "mtb"
  Midnasal,
  %- elif swab_type == "np"
  Nasopharyngeal,
  %- else
  \errmessage{Invalid swab_type "\VAR{swab_type}"}
  %- endif

  self-collected

\item [\textbf{Collection Barcode}] \VAR{qrcode|e}
\item [\textbf{Collection Date}] \VAR{collect_ts|e}
\item [\textbf{Report Date}]
  %- if status_code == "never-tested"
  Never tested
  %- else
  \VAR{result_ts|e}
  %- endif
\end{description}

\PageLine

\begin{center}
\Large
\textbf{RESEARCH RESULT - You are encouraged to share this result with your
  health care provider.}
\end{center}

\bigskip

\large \underline{\textbf{About This Test:}}

\begin{itemize}
\item

  COVID-19 is a \textbf{disease} that is caused by a \textbf{virus} called
  SARS-CoV-2. This research test looks for the virus in the sample that was
  collected from your nose to decide if you have COVID-19. The technical term
  for this kind of test is “reverse-transcription polymerase chain reaction,” or
  “RT-PCR” for short.

\end{itemize}

\bigskip

\large \underline{\textbf{Important Reminders}}

Before looking at your research test result, you should review the following:

\begin{itemize}
\item

  The \textbf{type of sample} that you collected and \textbf{the testing in the
  lab} are part of a \textbf{research study} and this report describes your
  \textbf{research test result.}

\item
%- if staff_observed == "false"
  You collected this sample and labeled it by yourself (or for someone else) and
  returned it to the lab via a delivery service. It was not collected, labeled,
  or transported by a licensed healthcare professional like a nurse or doctor.

  %- elif staff_observed == "true"
  You collected this sample by yourself (or for someone else) and it was labeled
  and returned to the lab via a courier or other type of non-healthcare
  professional. It was not collected or transported by a licensed healthcare
  professional like a nurse or doctor.

  %- else
  \errmessage{Unknown value: staff_observed = "\VAR{staff_observed}"}

  %- endif

  This may impact your test result, and a healthcare provider may want to test
  you again.

\item

  Being part of this \textbf{research study} does not take the place of your
  normal medical care. If you are feeling very sick or are concerned about your
  test result, you should contact a doctor or healthcare provider.

\end{itemize}

\bigskip
\PageLine

\large \underline{\textbf{Your SARS-CoV-2 (COVID-19) Research Test Result}}

%- if status_code == "negative"
Not Detected (Negative)\\
%- elif status_code == "positive"
Detected (positive)\\
%- elif status_code == "inconclusive"
Inconclusive\\
%- elif status_code == "never-tested"
Unable to be tested\\
%- else
\VAR{0/0} %# poor man's raise
%- endif

\PageLine
\bigskip

\large \underline{\textbf{Understanding Your Result}}

This list gives all of the results that your test may show and describes what
each of these results means. You can print out or email this \textbf{research
test result} to share with your doctor or healthcare provider

Regardless of your test result, you should stay home and avoid being around
other people if you are feeling sick. \textbf{If your symptoms continue, you are
very sick or worried about your illness, or have questions about your research
test result, you should contact a doctor or healthcare provider.}

\begin{itemize}


\item

  \textbf{If your test says ``Not Detected (Negative)''}, it means our lab did not find the
  virus that causes COVID-19 in the sample that you sent to us. No test is
  perfect, and there is a small chance that you could have the virus even though
  your test is negative.

\item

  \textbf{If your test says ``Detected (Positive)''}, it means our
  lab found the virus that causes COVID-19 in the sample that you sent us. You
  should stay home and avoid being around other people, especially if you are
  feeling sick.

  Your result will be shared with the Washington State and local health
  departments. You do \textbf{not} need to report these results yourself. The
  health department may call you to ask questions about your illness and people
  you may have been around while you were sick.

\item

  \textbf{If your test says ``Inconclusive''}, it means our lab is not certain
  of your test result. No test is perfect, and an inconclusive result happens
  when we find a small amount of SARS-CoV-2 in your sample, but not enough for
  us to say that you have COVID-19.

\item

  \textbf{If your test says ``Unable to be tested''}, it means there was either
  not enough material on your sample for us to test, or there was an error when
  we ran the test. Since we are not able to test your sample, we could not find
  whether or not you have the virus that causes COVID-19.

\end{itemize}

\bigskip

\large \underline{\textbf{Resources}}

If you would like more information on COVID-19, please review these resources
from public health partners:

\begin{itemize}
\item
  Public Health \textemdash Seattle \& King County\\
  \url{http://www.kingcounty.gov/covid}

\item

  Washington State Department of Health\\
  \url{http://www.doh.wa.gov/Emergencies/Coronavirus}

\item
  U.S. Centers for Disease Control and Prevention (CDC)\\
  \url{http://www.cdc.gov/coronavirus/2019-ncov/index.html}

\end{itemize}

If you would like to speak to someone on the phone, please call one of the
following:

\begin{itemize}

\item

  \textbf{Medical questions related to COVID-19:} Contact the King County Call
  Center between 8AM and 7PM at \textbf{206-477-3977}

\item

  \textbf{Non-medical questions about COVID-19, including compliance and
  business-related issues:} Contact the King County COVID-19 Business and
  Community Information Line, Monday – Friday between 8:30AM and 4:30PM at
  \textbf{206-296-1608}

\item

  \textbf{General questions about COVID-19 in Washington State:} Contact the
  Washington State Novel Coronavirus Call Center between 6AM and 10PM at
  \textbf{1-800-525-0127}

\end{itemize}

\bigskip

\large \underline{\textbf{Methods and Limitations}}

\textbf{Methods}\\
This assay is made to detect SARs-CoV-2 using reverse-transcription PCR. This
assay uses amplification of the virus segments ORf1B and S and also detects a
human RnaseP amplification control. This test was developed and its performance
characteristics determined by the University of Washington’s Northwest Genomics
Center. This laboratory is licensed by the WA Department of Health Medical
Testing Site as qualified to perform high complexity clinical laboratory
testing.

\textbf{Limitations}

\begin{itemize}

\item

  These results are provided as part of a research study.

\item
  %- if staff_observed == "false"
  You collected this sample and labeled it by yourself (or for someone else) and
  returned it to the lab via a delivery service. It was not collected, labeled,
  or transported by a healthcare professional like a nurse or doctor.

  %- elif staff_observed == "true"
  You collected this sample by yourself (or for someone else) and it was labeled
  and returned to the lab via a courier or other type of non-healthcare
  professional. It was not collected or transported by a licensed healthcare
  professional like a nurse or doctor.

  %- else
  \errmessage{Unknown value: staff_observed = "\VAR{staff_observed}"}

  %- endif

  This may impact your test result, and a healthcare provider may want to test
  you again.

\item

  Because this report is based on data from a research study, your results may
  not be able to diagnose COVID-19, the disease caused by an infection with
  SARS-CoV-2.

\end{itemize}

\end{document}
